\documentclass[journal]{IEEEtran}

\usepackage{iftex}
\ifPDFTeX
  \usepackage[utf8]{inputenc}
  \usepackage[T1]{fontenc}
\else
  \usepackage{fontspec}
  \setmainfont{texgyretermes}[Extension=.otf, UprightFont=*-regular, BoldFont=*-bold, ItalicFont=*-italic, BoldItalicFont=*-bolditalic]
\fi
\usepackage[ngerman]{babel}
\usepackage{amsmath, amssymb, amsfonts}
\usepackage{booktabs}
\usepackage{tabularx}
\usepackage{cite}
\usepackage{microtype}
\usepackage{url}
\usepackage{hyperref}

\hypersetup{
    colorlinks=true,
    linkcolor=black,
    citecolor=black,
    urlcolor=blue
}

\begin{document}

\title{Thermodynamische und techno-ökonomische Bewertung unterirdischer Naturzug-Schrägschächte zur passiven Kühlung und Energierückgewinnung in Hyperscale-KI-Rechenzentren}

\author{Chu Duc Kien
\thanks{Manuskript eingereicht am 11. September 2026. Korrespondenzadresse: \url{mailiekien@icloud.com}.}}

\markboth{IEEE Transactions on Sustainable Computing / Thermal Engineering,~Vol.~X, No.~X,~September~2026}%
{Chu Duc Kien: Unterirdische Naturzug-Schrägschächte für KI-Rechenzentren}

\maketitle

\begin{abstract}
Der exponentielle Anstieg von Trainingslasten künstlicher Intelligenz erfordert Rechenzentrumsleistungen im dreistelligen Megawatt-Bereich. Herkömmliche Kühlinfrastrukturen stoßen durch extremen Hilfsstrombedarf (Power Usage Effectiveness, PUE) und massiven Frischwasserverbrauch (Water Usage Effectiveness, WUE) an ökologische und infrastrukturelle Grenzen. Dieser Beitrag untersucht ein neuartiges Anlagenkonzept: die Kopplung eines flüssigkeitsgekühlten 100-MW-KI-Rechenzentrums mit einem bergmännisch aufgefahrenen Schrägschacht ($\Delta z = 1.500\,\text{m}$) zur Realisierung eines passiven Naturzug-Kühlsystems mit integrierter Energierückgewinnung. 

Thermodynamische Berechnungen zeigen, dass bei einer Kühlwassertemperatur von $60\,^\circ\text{C}$ ein kontinuierlicher Auftrieb mit einem Massenstrom von rund $4.480\,\text{kg/s}$ induziert wird. Neben der vollständigen Eliminierung des mechanischen Kühlstrombedarfs ($\sim 4{,}0\,\text{MW}_{\text{el}}$) und des jährlichen Wasserbedarfs ($\sim 1{,}58 \cdot 10^6\,\text{m}^3$) ermöglicht eine Basisturbine die Rückspeisung von rund $1{,}4\,\text{MW}_{\text{el}}$, wodurch der rechnerische PUE auf etwa $1{,}016$ sinkt. Eine techno-ökonomische Lebenszyklusanalyse ergibt jährliche Betriebskosteneinsparungen von rund $9{,}4\,\text{Mio. Euro}$; bei einer Netto-Mehrinvestition von $198\,\text{Mio. Euro}$ und einem Kalkulationszins von $6\,\%$ bleibt der Kapitalwert über 50 Jahre im Basisszenario jedoch negativ ($\approx -49\,\text{Mio. Euro}$). Ein positiver Kapitalwert stellt sich erst bei deutlich höheren Strom- oder Wasserpreisen ein. Das Konzept ist damit thermodynamisch tragfähig, wirtschaftlich aber standort- und preisabhängig.
\end{abstract}

\begin{IEEEkeywords}
Rechenzentrumskühlung, Direct Liquid Cooling, Naturzug, Schrägschacht, Energierückgewinnung, PUE, Thermosiphon, Techno-Ökonomie.
\end{IEEEkeywords}

\section{Einleitung}
\IEEEPARstart{D}{ie} fortschreitende Skalierung tiefer künstlicher neuronaler Netze, insbesondere autoregressiver Foundation-Modelle, induziert einen beispiellosen Anstieg der Leistungsdichte in modernen Rechenzentren. Während klassische Cloud-Infrastrukturen historisch auf Wärmelasten von $5\,\text{bis}\,15\,\text{kW}$ pro Rack ausgelegt waren, überschreiten dichte GPU-Cluster der aktuellen Generation regelmäßig Grenzwerte von $40\,\text{bis}\,120\,\text{kW}$ pro Schrankeinheit \cite{ashrae2023}. Rein luftgeführte Kühlsysteme stoßen hierbei an physikalische Grenzen: Die notwendigen Volumenströme verursachen prohibitive Druckverluste und einen stark ansteigenden Hilfsenergiebedarf für die interne Ventilation.

Die Industrie reagiert hierauf mit dem flächendeckenden Übergang zu Direct Liquid Cooling (DLC). Da flüssige Wärmeträger eine um Größenordnungen höhere volumetrische Wärmekapazität als Luft aufweisen, erlaubt die Direktkühlung der Prozessorkerne eine Erhöhung der Rücklauftemperaturen auf ein Niveau von $55\,\text{bis}\,65\,^\circ\text{C}$. Trotz dieser thermodynamischen Verbesserung bleibt die Abführung thermischer Megawattströme problematisch:
\begin{itemize}
    \item \textbf{Trockenkühlung:} Lamellen-Rückkühler erfordern bei hohen sommerlichen Umgebungstemperaturen erhebliche elektrische Ventilatorleistungen oder den Einsatz von Kompressionskältemaschinen (Chiller), welche den PUE drastisch verschlechtern.
    \item \textbf{Verdunstungskühlung:} Die Ausnutzung der Verdampfungsenthalpie reduziert zwar die elektrische Hilfsenergie, erzeugt jedoch einen extremen Wasserfußabdruck. Ein 100-MW-Rechenzentrum verbraucht unter Volllast jährlich zwischen $1{,}5$ und $2{,}0\,\text{Mrd. Litern}$ Frischwasser. In ariden Gebieten führt dies zunehmend zu regulatorischen Baustopps.
\end{itemize}

Vor diesem Hintergrund analysiert die vorliegende Arbeit die thermodynamische und techno-ökonomische Machbarkeit eines unterirdischen Naturzug-Schrägschachts als passives Kühlsystem für ein 100-MW-KI-Rechenzentrum im alpinen Gelände.

\section{Systemarchitektur und Prozesskopplung}
Das Anlagenkonzept gliedert sich in drei funktionale Einheiten: die Serverkaverne am Fuße des Bergmassivs, die Wärmeübertragungs- und Turbinenstation sowie den bergmännisch aufgefahrenen Schrägschacht.

\subsection{Kühlkreislauf-Topologie}
Die Server-Hardware wird über geschlossene Primärkreisläufe mit deionisiertem Wasser gekühlt. Bei einer Eintrittstemperatur von $35\,^\circ\text{C}$ und einer Austrittstemperatur von $60\,^\circ\text{C}$ transportiert das Fluid rund $90\,\text{MW}_{\text{th}}$ thermische Verlustleistung zu einem Sekundärwärmetauscher. 

\subsection{Aerodynamische Schachtgeometrie}
Zur Gewährleistung einer redundanten Betriebsführung wird der Schacht als Doppelröhrensystem mit je $D_{\text{hyd}} = 13{,}4\,\text{m}$ ausgeführt. Bei einer geodätischen Höhendifferenz von $H = 1.500\,\text{m}$ und einem mittleren Neigungswinkel von $\alpha = 30^\circ$ beträgt die Auffahrlänge $L = 3.000\,\text{m}$. Der Gesamteinströmquerschnitt am Schachtfuß beträgt $A_{\text{ges}} \approx 282\,\text{m}^2$; der Durchmesser ist so gewählt, dass sich beim Auslegungsmassenstrom eine mittlere Strömungsgeschwindigkeit von rund $15\,\text{m/s}$ einstellt.

\section{Thermodynamische Modellierung}

\subsection{Erhaltungsgleichungen}
Für die stationäre, eindimensionale Strömung entlang der Schachtkoordinate $s$ gilt unter Vernachlässigung von Massenquellen:
\begin{equation}
\dot{m} = \rho(z) \cdot A(z) \cdot v(z) = \text{const.}
\end{equation}

Der dem Luftstrom über die Wärmetauscher zugeführte Wärmestrom $\dot{Q}_{\text{th}}$ bestimmt den Temperaturanstieg $\Delta T = T_i - T_a$:
\begin{equation}
\dot{Q}_{\text{th}} = \dot{m} \cdot c_p \cdot (T_i - T_a)
\end{equation}
Hierbei bezeichnen $c_p$ die spezifische Wärmekapazität von Luft ($1.005\,\text{J/(kg}\cdot\text{K)}$), $T_a$ die atmosphärische Eintrittstemperatur und $T_i$ die Mischtemperatur im Schacht nach Wärmeaufnahme.

\subsection{Antriebsdruckdifferenz (Kamineffekt)}
Der statische Gesamtauftrieb $\Delta p_{\text{stat}}$ ergibt sich durch die hydrostatische Dichtedifferenz zwischen der ungestörten Atmosphärensäule $\rho_a(z)$ und der erwärmten Luftsäule $\rho_i(z)$ über die Schachtachse \cite{schlaich1995}:
\begin{equation}
\Delta p_{\text{stat}} = g \int_0^H \left[ \rho_a(z) - \rho_i(z) \right] dz
\end{equation}

Unter Verwendung der idealen Gasgleichung $\rho = p / (R_s T)$ und Linearisierung der Temperaturprofile über $H$ lässt sich der Auftrieb näherungsweise beschreiben durch:
\begin{equation}
\Delta p_{\text{stat}} \approx \bar{\rho}_a \cdot g \cdot H \cdot \left( \frac{T_i - T_a}{T_i} \right)
\end{equation}
Mit $\bar{\rho}_a \approx 1{,}10\,\text{kg/m}^3$, $H = 1.500\,\text{m}$, $T_a = 293{,}15\,\text{K}$ ($20\,^\circ\text{C}$) und $\Delta T = 20\,\text{K}$ ergibt sich:
\begin{equation}
\Delta p_{\text{stat}} \approx 1{,}10 \cdot 9{,}81 \cdot 1.500 \cdot \left(\frac{20}{313{,}15}\right) \approx 1.033\,\text{Pa}
\end{equation}

\subsection{Druckverlustbilanz und Turbinenleistung}
Die Fließgeschwindigkeit stellt sich im dynamischen Gleichgewicht zwischen Antriebsdruck und dissipativen Verlusten ein:
\begin{equation}
\Delta p_{\text{stat}} = \Delta p_{\text{WT}} + \Delta p_{\text{fric}} + \Delta p_{\text{form}} + \Delta p_{\text{turb}}
\end{equation}

Die Wandreibung im Schacht wird über den Darcy-Weisbach-Ansatz quantifiziert:
\begin{equation}
\Delta p_{\text{fric}} = \lambda \cdot \frac{L}{D_{\text{hyd}}} \cdot \frac{\rho_i v^2}{2}
\end{equation}
Für geglätteten Spritzbeton ($\varepsilon \approx 1{,}5\,\text{mm}$) und voll ausgebildete turbulente Strömung ($\text{Re} > 10^7$) wird $\lambda \approx 0{,}013$ gesetzt. Mit dem Massenstrom $\dot{m} \approx 4.478\,\text{kg/s}$ aus Gl.~(2), der Dichte der erwärmten Schachtluft $\rho_i \approx 1{,}05\,\text{kg/m}^3$ und $A_{\text{ges}} \approx 282\,\text{m}^2$ ergibt sich ein volumetrischer Durchsatz $\dot{V} = \dot{m}/\rho_i \approx 4.264\,\text{m}^3/\text{s}$ und eine mittlere Geschwindigkeit $v \approx 15{,}1\,\text{m/s}$; daraus folgt $\Delta p_{\text{fric}} \approx 349\,\text{Pa}$. Unter Einbeziehung des Wärmetauscherdruckverlusts ($\Delta p_{\text{WT}} \approx 180\,\text{Pa}$) und von Einlaufformverlusten ($\Delta p_{\text{form}} \approx 80\,\text{Pa}$) verbleibt für die Turbine ein Druckgefälle von
\begin{equation}
\Delta p_{\text{turb}} \approx 1.033 - 180 - 349 - 80 \approx 424\,\text{Pa}.
\end{equation}

Die über Axialturbinen am Schachtfuß rekuperierte elektrische Leistung $P_{\text{el}}$ berechnet sich mit dem Wirkungsgrad $\eta_{\text{turb}} = 0{,}78$ (Turbine und Generator) zu:
\begin{equation}
P_{\text{el}} = \eta_{\text{turb}} \cdot \dot{V} \cdot \Delta p_{\text{turb}} \approx 0{,}78 \cdot 4.264 \cdot 424 \approx 1{,}41\,\text{MW}_{\text{el}}
\end{equation}
Die Turbine kann somit rund $41\,\%$ des statischen Auftriebs nutzen; der übrige Anteil wird durch Wärmetauscher, Wandreibung und Formverluste dissipiert. Ein größerer Turbinenanteil wäre nur durch einen geringeren Reibungsverlust (größerer Schachtquerschnitt, geringere Geschwindigkeit) erreichbar, was jedoch die Vortriebskosten erhöht.

\subsection{Theoretischer Kaminwirkungsgrad}
Das thermodynamische Limit der Wandlung von Niedertemperaturwärme in Strömungsarbeit über einen Auftriebskamin ist unabhängig von der Temperaturerhöhung $\Delta T$ und skaliert ausschließlich mit der Bauhöhe $H$ \cite{haaf1983}:
\begin{equation}
\eta_{\text{kamin}} = \frac{g \cdot H}{c_p \cdot T_a} = \frac{9{,}81 \cdot 1.500}{1.005 \cdot 293{,}15} \approx 4{,}99\,\%
\end{equation}
Bezogen auf $\dot{Q}_{\text{th}} = 90\,\text{MW}_{\text{th}}$ entspricht dies einer idealen Strömungsleistung von rund $4{,}5\,\text{MW}$. Der aus der Druckverlustbilanz abgeleitete Wert von $1{,}41\,\text{MW}_{\text{el}}$ entspricht einem Gesamtwirkungsgrad von etwa $1{,}6\,\%$ bezogen auf die zugeführte Wärme, also rund einem Drittel des theoretischen Limits. Dieser Wert ist konsistent mit den bei Aufwindkraftwerken beobachteten Gesamtwirkungsgraden \cite{haaf1983, ho2021}.

\begin{table}[htbp]
\caption{Systemparameter des 100-MW-Referenzsystems}
\label{tab:system_params}
\centering
\begin{tabularx}{\columnwidth}{l l r l}
\toprule
\textbf{Parameter} & \textbf{Symbol} & \textbf{Wert} & \textbf{Einheit} \\
\midrule
IT-Leistungsaufnahme & $P_{\text{IT}}$ & 100,0 & $\text{MW}_{\text{el}}$ \\
Thermische Verlustleistung & $\dot{Q}_{\text{th}}$ & 90,0 & $\text{MW}_{\text{th}}$ \\
Geodätische Schafthöhe & $H$ & 1.500 & $\text{m}$ \\
Schachtlänge ($\alpha = 30^\circ$) & $L$ & 3.000 & $\text{m}$ \\
Hydraulischer Durchmesser (je Röhre) & $D_{\text{hyd}}$ & 13,4 & $\text{m}$ \\
Gesamtquerschnitt & $A_{\text{ges}}$ & 282 & $\text{m}^2$ \\
Temperaturhub Zuluft/Schacht & $\Delta T$ & 20,0 & $\text{K}$ \\
Luftmassenstrom & $\dot{m}$ & 4.478 & $\text{kg/s}$ \\
Volumenstrom & $\dot{V}$ & 4.264 & $\text{m}^3/\text{s}$ \\
Mittlere Strömungsgeschwindigkeit & $v$ & 15,1 & $\text{m/s}$ \\
Statischer Auftrieb & $\Delta p_{\text{stat}}$ & 1.033 & $\text{Pa}$ \\
Turbinendruckgefälle & $\Delta p_{\text{turb}}$ & 424 & $\text{Pa}$ \\
Generierte Turbinenleistung & $P_{\text{gen}}$ & 1,41 & $\text{MW}_{\text{el}}$ \\
Vermiedene Lüfterhilfsleistung & $P_{\text{verm}}$ & 4,00 & $\text{MW}_{\text{el}}$ \\
\bottomrule
\end{tabularx}
\end{table}

\section{Energetische und ökologische Performance}

\subsection{Modifizierte PUE-Metrik}
Die Standardmetrik des Konsortiums The Green Grid definiert die Power Usage Effectiveness als:
\begin{equation}
\text{PUE} = \frac{P_{\text{gesamt}}}{P_{\text{IT}}} = \frac{P_{\text{IT}} + P_{\text{Klima}} + P_{\text{Verlust}}}{P_{\text{IT}}}
\end{equation}

Da der Hilfsenergieaufwand für die Kühlung $P_{\text{Klima}} = 0$ beträgt und das Kühlsystem über die Kaminturbinen Nettoenergie $P_{\text{gen}}$ bereitstellt, modifiziert sich die Energiebilanz:
\begin{equation}
\text{PUE}_{\text{mod}} = \frac{P_{\text{IT}} + P_{\text{Verlust}} - P_{\text{gen}}}{P_{\text{IT}}}
\end{equation}
Unter Berücksichtigung interner Wandlungs- und Verteilverluste ($P_{\text{Verlust}} \approx 0{,}03 \cdot P_{\text{IT}}$) sinkt der PUE auf:
\begin{equation}
\text{PUE}_{\text{mod}} \approx \frac{100\,\text{MW} + 3{,}0\,\text{MW} - 1{,}41\,\text{MW}}{100\,\text{MW}} \approx 1{,}016
\end{equation}
Auch ohne Anrechnung der Turbinenleistung liegt der PUE bei $1{,}03$ und damit deutlich unter typischen Werten luft- oder verdunstungsgekühlter Hyperscale-Anlagen ($1{,}15$--$1{,}25$).

\subsection{Wasserbedarfsbilanzierung (WUE)}
Herkömmliche Kühlsysteme mit evaporativer Unterstützung verbrauchen etwa $1{,}8\,\text{L/kWh}_{\text{IT}}$. Bei kontinuierlicher Volllast ($\tau = 8.760\,\text{h/a}$) entspricht dies:
\begin{equation}
V_{\text{H2O}} = 100.000\,\text{kW} \cdot 8.760\,\text{h/a} \cdot 1{,}8\,\text{L/kWh} \approx 1{,}58 \cdot 10^6\,\text{m}^3/\text{a}
\end{equation}
Durch den Einsatz des Schrägschachts sinkt die Water Usage Effectiveness des Rechenzentrums auf $\text{WUE} = 0{,}00\,\text{L/kWh}$.

\section{Geotechnische Auslegung und Tunnelbau}
Die Realisierung eines Schrägschachts dieser Dimensionierung lehnt sich an Entwürfe moderner Hochdruck-Pumpspeicherkraftwerke an (z.\,B. Nant de Drance, Linthal 2015).

\subsection{Vortriebsverfahren}
Zur Erstellung der Zwillingsröhre kommt primär eine Kombination aus Raise-Boring und maschinellem Schrägvortrieb mittels Gripper-Tunnelbohrmaschine (TBM) zum Einsatz. Raise-Boring-Systeme ermöglichen das pilotierte Bohren von unten nach oben mit anschließendem Rückzugserweitern, wodurch anfallendes Ausbruchmaterial gravitationsbedingt direkt am Fuße des Schachts abgezogen werden kann.

\subsection{Gebirgsmechanische Stabilität}
Im Granit- und Gneisgebirge erfordert die kontinuierliche thermische Beaufschlagung der Schachtlaibung ($T_{\text{wand}} \approx 40\,^\circ\text{C}$) eine strukturmechanische Sicherung gegen hitzeinduzierte Gebirgsspannungen. Der Ausbau erfolgt über drainierten Spritzbeton mit Systemankerung zur Vermeidung von hydrostatischem Kluftwasserdruck.

\section{Techno-ökonomische Bewertung}

\subsection{Differenzinvestitionskosten ($\Delta\text{CapEx}$)}
Die Baukosten des Schrägschachtsystems werden mit dem Einsparungspotenzial entfallender konventioneller Kühlanlagen saldiert:

\begin{itemize}
    \item \textbf{Schachtauffahrung:} $2 \times 3.000\,\text{m}$ Vortrieb ($32.000\,\text{EUR/m}$): $192{,}0\,\text{Mio.\,EUR}$. Der Meterpreis ist ein Richtwert am unteren Rand alpiner Referenzprojekte; die Kostenabhängigkeit vom Querschnitt wird über das Szenario \glqq{}Baukosten $+25\,\%$\grqq{} in Tabelle~\ref{tab:sensitivity} abgedeckt.
    \item \textbf{Felsausbau und Glättung:} Spritzbetonschale und Felsanker: $18{,}0\,\text{Mio.\,EUR}$.
    \item \textbf{Portalanlagen \& Geotechnik:} Diffusoren, Schutzbauten: $12{,}0\,\text{Mio.\,EUR}$.
    \item \textbf{Wärmetauscher und Turbinen:} $14{,}5\,\text{Mio.\,EUR}$.
    \item \textbf{Brutto-Investition Schacht ($\text{CapEx}_{\text{sys}}$):} $236{,}5\,\text{Mio.\,EUR}$.
\end{itemize}

Dem gegenüber steht der Wegfall klassischer Chiller, Kühltürme und Wasseraufbereitungsanlagen ($\text{CapEx}_{\text{konv}} \approx 38{,}5\,\text{Mio.\,EUR}$), was eine Netto-Mehrinvestition von $\Delta\text{CapEx} = 198{,}0\,\text{Mio.\,EUR}$ ergibt.

\subsection{Betriebskosteneinsparung ($\Delta\text{OPEX}$)}
Bei Standardpreisen von $k_{\text{el}} = 0{,}10\,\text{EUR/kWh}$ und $k_{\text{w}} = 2{,}50\,\text{EUR/m}^3$ errechnen sich die jährlichen OPEX-Vorteile zu:
\begin{align}
\Delta C_{\text{p-klima}} &= 4.000\,\text{kW} \cdot 8.760\,\text{h} \cdot 0{,}10\,\text{EUR/kWh} = 3{,}50\,\text{Mio. EUR/a} \\
R_{\text{gen}} &= 1.410\,\text{kW} \cdot 8.760\,\text{h} \cdot 0{,}10\,\text{EUR/kWh} = 1{,}24\,\text{Mio. EUR/a} \\
\Delta C_{\text{wasser}} &= 1.576.800\,\text{m}^3 \cdot 2{,}50\,\text{EUR/m}^3 = 3{,}94\,\text{Mio. EUR/a} \\
\Delta C_{\text{maint}} &\approx 0{,}75\,\text{Mio. EUR/a}
\end{align}
Daraus folgt eine jährliche Gesamtentlastung von $\mathbf{\Delta\text{OPEX} \approx 9{,}43\,\text{Mio. EUR/a}}$. Die vermiedene Kühlhilfsenergie und der vermiedene Wasserbezug tragen zusammen rund $79\,\%$ bei; die Turbinenerlöse sind mit $13\,\%$ ein Nebeneffekt, nicht der wirtschaftliche Kern des Konzepts.

\subsection{Kapitalwert und Sensitivitätsanalyse}
Der Kapitalwert (Net Present Value, NPV) berechnet sich bei einem Kalkulationszinssatz $r = 6\,\%$ über $N$ Jahre:
\begin{equation}
\text{NPV} = -\Delta\text{CapEx} + \sum_{t=1}^N \frac{\Delta\text{OPEX}_t}{(1 + r)^t}
\end{equation}
Bei konstantem $\Delta\text{OPEX}$ entspricht die Summe dem Rentenbarwertfaktor $AF(r,N) = \left[1-(1+r)^{-N}\right]/r$; für $r = 6\,\%$ und $N = 50$ ergibt sich $AF \approx 15{,}76$. Im Basisszenario folgt daraus
\begin{equation}
\text{NPV}_{50} \approx -198{,}0 + 9{,}43 \cdot 15{,}76 \approx -49{,}3\,\text{Mio. EUR}.
\end{equation}
Die Nullgrenze wird bei $\Delta\text{OPEX} \approx 12{,}6\,\text{Mio. EUR/a}$ erreicht; die undiskontierte Amortisationszeit beträgt rund 21 Jahre.

\begin{table}[htbp]
\caption{Sensitivitätsanalyse des Kapitalwerts ($r = 6\,\%$)}
\label{tab:sensitivity}
\centering
\footnotesize\setlength{\tabcolsep}{4pt}
\begin{tabular}{@{}l r r r r@{}}
\toprule
\textbf{Szenario} & $k_{\text{el}}$ & $k_{\text{w}}$ & $\Delta\text{OPEX}$ & $\text{NPV}_{50}$ \\
 & [EUR/kWh] & [EUR/m$^3$] & [Mio.\,EUR/a] & [Mio.\,EUR] \\
\midrule
Basisszenario & 0,10 & 2,50 & 9,43 & $-49{,}3$ \\
Hohe Stromkosten & 0,16 & 2,50 & 12,27 & $-4{,}5$ \\
Wasserknappheit & 0,10 & 5,00 & 13,37 & $+12{,}8$ \\
Baukosten $+25\,\%$ & 0,10 & 2,50 & 9,43 & $-98{,}8$ \\
Kombiniert (Future) & 0,15 & 4,00 & 14,17 & $+25{,}3$ \\
\bottomrule
\end{tabular}
\end{table}

Über einen für alpine Infrastrukturbauten üblichen Betrachtungshorizont von $N = 50\,\text{Jahren}$ ist der Kapitalwert im Basisszenario negativ. Positiv wird er erst unter Bedingungen, die in wasserarmen Regionen oder bei anhaltend hohen Strompreisen eintreten können (Wasserpreis $\geq 5\,\text{EUR/m}^3$ oder kombiniert $0{,}15\,\text{EUR/kWh}$ und $4\,\text{EUR/m}^3$). Eine Kostenüberschreitung im Tunnelbau um $25\,\%$, wie sie bei Großprojekten dieser Art regelmäßig auftritt, verschlechtert den Kapitalwert um rund $50\,\text{Mio. Euro}$. Die Wirtschaftlichkeit des Konzepts hängt somit weniger von der Thermodynamik als von Standortfaktoren, Energie- und Wasserpreisentwicklung sowie der Kostendisziplin im Vortrieb ab. Nicht monetär bewertet sind der Wegfall regulatorischer Wasserrisiken und der CO$_2$-Fußabdruck der vermiedenen Hilfsenergie.

\section{Diskussion, Restriktionen und Ausblick}

\subsection{Betriebssicherheit bei Inversion und Vereisung}
Kritische Betriebszustände entstehen bei meteorologischen Inversionswetterlagen, wenn der vertikale Dichtegradient der freien Atmosphäre temporär einbricht. Um einen Stillstand oder eine Strömungsumkehr (*Backdraft*) zu unterbinden, müssen die Turbinenaggregate als umkehrbare Axialventilatoren ausgeführt werden. 

Im Winterhalbjahr führt die Entspannung der aufsteigenden Luftmasse am Gipfelaustritt zu Kondensations- und Reifbildung. Eine gezielte Beimischung warmer Schachtluft im Diffusorbereich verhindert Vereisungen an den Schutzgittern.

\subsection{Standortfaktoren}
Die Realisierung erfordert spezifische topografische Rahmenbedingungen (Hangneigung $> 25^\circ$, Granit-/Gneisformationen) sowie eine adäquate Netzanbindung. Ideale Standorte bieten Alpentäler in der Nähe bestehender Speicherkraftwerke, da hier bereits redundante Hochspannungstrassen ($110\text{--}220\,\text{kV}$) existieren.

\section{Fazit}
Die Integration eines unterirdischen Naturzug-Schrägschachts in ein 100-MW-KI-Rechenzentrum adressiert die thermischen und ökologischen Kernprobleme moderner Hyperscale-Infrastrukturen: Der Kühlbetrieb kommt ohne mechanische Hilfsenergie und ohne Frischwasser aus, der PUE sinkt auf etwa $1{,}016$ (WUE $= 0{,}0\,\text{L/kWh}$), und eine Basisturbine speist rund $1{,}4\,\text{MW}_{\text{el}}$ zurück. Thermodynamisch ist das Konzept damit robust.

Wirtschaftlich ist es im Basisszenario nicht tragfähig: Den jährlichen Betriebskosteneinsparungen von rund $9{,}4\,\text{Mio. Euro}$ steht eine Netto-Mehrinvestition von $198\,\text{Mio. Euro}$ gegenüber, was bei $6\,\%$ Kalkulationszins über 50 Jahre einen Kapitalwert von etwa $-49\,\text{Mio. Euro}$ ergibt. Ein positiver Kapitalwert setzt hohe Wasser- oder Strompreise voraus, wie sie in wasserarmen Regionen bereits heute auftreten oder mittelfristig erwartet werden. Das Konzept ist daher weniger eine allgemeine Blaupause als eine standortspezifische Option für Regionen, in denen Wasserverfügbarkeit oder Energiepreise den Bau konventioneller Kühlsysteme einschränken. Weiterführende Arbeiten sollten die Reibungsverluste über CFD-Simulationen absichern, den Turbinenwirkungsgrad bei Druckdifferenzen unter $500\,\text{Pa}$ experimentell belegen und die Vortriebskosten anhand konkreter Standorte präzisieren.

\section*{Revisionsvermerk}
Diese Fassung korrigiert die Erstfassung des Preprints in drei Punkten: (1) Die in der Erstfassung genannte Turbinenleistung von $2{,}25\,\text{MW}_{\text{el}}$ ließ sich aus der dort angegebenen Druckverlustbilanz nicht reproduzieren (korrekt: $\approx 1{,}4\,\text{MW}_{\text{el}}$); (2) die Schachtgeometrie ($D_{\text{hyd}}$, $A_{\text{ges}}$, $v$) war intern inkonsistent und wurde auf $v \approx 15\,\text{m/s}$ vereinheitlicht; (3) der Kapitalwert von $+162{,}3\,\text{Mio. Euro}$ im Basisszenario war aus der angegebenen NPV-Formel und den angegebenen Eingangsgrößen nicht herleitbar (korrekt: $\approx -49\,\text{Mio. Euro}$). PUE, OPEX, Sensitivitätsanalyse, Abstract und Fazit wurden entsprechend angepasst. Die Fehler wurden durch eine unabhängige Nachrechnung mit KI-Unterstützung aufgedeckt.

\begin{thebibliography}{1}

\bibitem{ashrae2023}
ASHRAE Technical Committee 9.9, \emph{Liquid Cooling Guidelines for Datacom Equipment Centers}, 3rd~ed., Atlanta, GA: ASHRAE, 2023.

\bibitem{schlaich1995}
J.~Schlaich, \emph{The Solar Chimney: Electricity from the Sun}, Stuttgart, Deutschland: Edition Axel Menges, 1995.

\bibitem{haaf1983}
W.~Haaf, K.~Friedrich, G.~Mayr, und J.~Schlaich, \glqq{}Solar Chimneys: Part I: Principle and Construction of the Manzanares Pilot Plant\grqq{}, \emph{International Journal of Solar Energy}, Bd.~2, Nr.~1, S.~3--20, 1983.

\bibitem{thegreengrid2012}
The Green Grid, \emph{PUE: A Comprehensive Examination of the Metric}, White Paper \#49, Beaverton, OR: The Green Grid, 2012.

\bibitem{ho2021}
C.~K.~Ho, \glqq{}A review of solar updraft towers: Modeling, design, and performance\grqq{}, \emph{Renewable and Sustainable Energy Reviews}, Bd.~146, Art.-Nr.~111162, 2021.

\bibitem{ita2022}
International Tunnelling and Underground Space Association (ITA), \emph{Guidelines for the Design of Inclined Shafts and Pressure Tunnels in Rock}, ITA Working Group 2, Report No.~24, Genf, 2022.

\end{thebibliography}

\end{document}
